% Migrated from proposal/proposal.tex (5 October 2026) on 14 October 2026.
% Written out on 12 November 2026 from the commented proposal text; the facts were
% checked against the sources cited.
\chapter{Background}
\label{ch:background}

This chapter introduces the apps the study analyses, the legal rules on what
their providers must tell users, the distinction between a privacy violation and
a privacy harm, and the protective techniques whose claims the study examines.

\section{Health Tracking Apps}
\label{sec:bg-trackers}

In Germany, 73\,\% of smartphone users use apps to support their fitness,
nutrition or general wellbeing. In a representative survey of 1{,}145 people
aged 16 and over, 43\,\% used a central health app such as Apple Health, Google
Fit or Health Connect, 37\,\% used sport apps that record activities such as
running or cycling, and 30\,\% used apps that record body data such as heart
rate, sleep or breathing~\parencite{bitkom2025health}.

Health trackers combine information that users enter, such as symptoms, meals or
cycle dates, with sensor readings and usage records. Together, these data can
reveal more than users directly provide. A survey of 236 studies of wearable
activity trackers reports that their data can be mined to infer private
information such as the user's personality traits~\parencite{niksirat2024wearable}.
The study covers such apps on the German market, including the companion apps
of wearables (Section~\ref{sec:rq-scope}).

\section{Legal Framework}
\label{sec:bg-law}

\paragraph{Health data.} Data concerning health and data concerning a person's
sex life are special categories of personal data under Article~9(1) of the
General Data Protection Regulation (GDPR). Data concerning health are personal
data that reveal information about a person's health status (Article~4(15) and
Recital~35)~\parencite{gdpr2016}. The Court of Justice of the European Union
reads these categories widely. Data that can reveal sensitive information
\enquote{by means of an intellectual operation involving comparison or
deduction} fall within Article~9(1)~\parencite[paras.~120, 125]{cjeu2022ot},
and the data customers enter when ordering pharmacy-only
medicines online are data concerning health even where it is only probable that
the medicines are meant for them~\parencite[para.~90]{cjeu2024lindenapotheke}.
Activity, sleep or cycle records can therefore be health data when they allow
such conclusions.

\paragraph{Transparency.} Article~5(1)(a) requires that personal data be
processed lawfully, fairly and transparently, and Article~12 requires that
information be given in a concise, transparent, intelligible and easily
accessible form, using clear and plain language. Articles~13 and~14 list what a
controller must tell data subjects, among them the purposes of processing, the
recipients and the data subjects' rights~\parencite{gdpr2016}. Guidance of the
Article~29 Working Party, endorsed by the European Data Protection Board, allows
this information to be layered across a privacy statement and notices in the
interface~\parencite{wp292018transparency}. The policy and the interface are
therefore both places where a provider informs users. Misleading commercial
claims, including claims about privacy, are also subject to consumer-protection
law~\parencite{ucpd2005}. \highlight[orange!25]{The study does not assume that
all communication about harms is voluntary. It asks what providers say about
possible consequences alongside their data-practice disclosures.}

\paragraph{Wellness applications in the European Health Data Space.} The
Regulation on the European Health Data Space defines a \emph{wellness
application} as software intended by its manufacturer to be used by a natural
person for processing electronic health data, specifically for providing
information on the health of natural persons or the delivery of care for
purposes other than the provision of healthcare
(Article~2(2)(ab))~\parencite{ehds2025}. A wellness application needs a label
only if its manufacturer claims interoperability with electronic health record
systems (Article~47), and data from wellness applications are among the
minimum categories of electronic health data made available for secondary use
(Article~51(1)(i)). The Regulation applies from 26~March 2027, with later dates
for some provisions, including Article~47 (Article~105)~\parencite{ehds2025}.

\section{Privacy Violations and Privacy Harms}
\label{sec:bg-harms}

A privacy violation and a privacy harm are related but distinct. Calo separates
subjective harm, the perception of unwanted observation, from objective harm,
the unanticipated or coerced use of information against a
person~\parencite{calo2011boundaries}. Citron and Solove describe the harms that
follow from privacy violations, among them physical, economic, reputational,
psychological, autonomy, discrimination and relationship
harms~\parencite{citron2021harms}. Knabenhans et al.\ use \emph{harm trees} to
study the European Digital Identity
Framework~\parencite[\S4.1]{knabenhans2026harms}. A tree links system
components, such as information disclosed to another party, to capabilities,
such as tracking or inference, and then to possible harms. Components may arise
from the core functionality, its digital context or implementation choices.
Their analysis distinguishes pathways to harm that are inherent to the core
functionality from pathways that privacy-enhancing technologies can prevent, and
shows that these technologies reduce information flows but fall short of
mitigating the harms of that system~\parencite{knabenhans2026harms}. This study
adapts the method to health trackers (Section~\ref{sec:framework}). It
distinguishes an app's intended prediction for its user from an adversary's
ability to obtain or use that prediction.

\section{Privacy-Enhancing Technologies and Protection Claims}
\label{sec:bg-pets}

Providers describe measures that protect users' data, such as encryption,
pseudonymisation or de-identification, aggregation and processing on the device.
Martin et al.\ examine three principles behind the privacy-enhancing techniques
of large technology companies: limiting third-party access to personal data,
using inferences instead of raw data, and keeping data on the user's device.
They argue that these principles fall short when the technique does not address
the privacy-violating behaviour, and report a gap between the principles and
users' privacy expectations~\parencite{martin2025cookies}. A measure can thus do
what it promises and still leave a harm possible. The relevant question for this
study is which protection a provider claims, against which actor and under which
conditions, and which pathways to harm the measure interrupts.
